% !TeX root = ../main.tex

\chapter{資料特徵與真實需求還原}\label{chap:Data & Latent Demand Reconstruction}

\section{\textcolor{red}{資料來源與時空特徵}}\label{sec:data_sources}

\textcolor{red}{（這 2.1 的數據跟數值全都要等結束後重新校正，先寫而已，然後真實的處理流程全都要記錄下來，看要怎樣刪減。）}

\subsection{多源異質資料集架構}\label{subsec:data_sources_overview}

為精確捕捉國立臺灣大學校總區及周邊生活圈之共享單車供需動態，本研究整合四類來自政府開放資料平臺與學術機構之異質實證資料集，涵蓋微觀票證交易流水、高頻站點車位快照、局屬氣象觀測以及學術與行政日曆。各資料集之來源、觀測期間與樣本規模彙整於表~\ref{tab:data_sources_summary}，其具體結構與時空特徵說明如下：

\begin{enumerate}
    \item \textbf{臺北市公共自行車 2.0 電子票證租借交易紀錄}：
    取自臺北市資料大平臺按月釋出之 YouBike 2.0 票證刷卡資料，觀測期間涵蓋 2024 年 1 月至 2025 年 12 月共 24 個月份檔案，全市原始交易流水總計達 $133{,}966{,}680$ 筆（其中 2024 年為 $55{,}881{,}184$ 筆，2025 年擴增至 $78{,}085{,}496$ 筆）。每筆紀錄記載單次完成扣款旅次之借車時間（\texttt{rent\_time}）、借車站點名稱（\texttt{rent\_station}）、還車時間（\texttt{return\_time}）、還車站點名稱（\texttt{return\_station}）、精確至秒之實際騎乘歷時（\texttt{rent}，格式為 \texttt{H+:MM:SS}）以及資料歸屬日期（\texttt{infodate}）。基於去識別化隱私保護規範，官方釋出之借還時間戳記將分與秒截斷為整點（\texttt{:00:00}），提供小時級之起訖站點流量（Origin-Destination, OD）資訊。此外，該資料集於 2024 年 11 月發生結構演進：2024 年 1 月至 10 月採 6 欄位格式（未區分車種）；自 2024 年 11 月起新增第 7 欄位車種標記（\texttt{bike\_type}），始明確區分「一般車」（YouBike 2.0）與「電輔車」（YouBike 2.0E）。

    \item \textbf{交通部 TDX 站點歷史車位即時快照}：
    取自交通部運輸資料流通服務平臺（TDX）之臺北市公共自行車站點剩餘數量歷史資料（資料標準結構為 \texttt{BikeAvailability}）。2024 全閏年共計 366 個每日原始 CSV 檔案，全市累計達 $111{,}802{,}383$ 筆狀態快照。每筆快照包含 10 個原始欄位：站點唯一識別碼 \texttt{StationUID}（具 \texttt{TPE} 前綴，如 \texttt{TPE500101001}）、9 碼站點代碼 \texttt{StationID}、營運狀態碼 \texttt{ServiceStatus}、服務型態 \texttt{ServiceType}（全數為 \texttt{YouBike2}）、帳面可借車數 \texttt{AvailableRentBikes}（即 $I_{i,\tau}$）、可還空柱數 \texttt{AvailableReturnBikes}（即 $O_{i,\tau}$）、車柱設備狀態更新時間 \texttt{SrcUpdateTime}、平臺伺服器接收時間 \texttt{UpdateTime}，以及一般車數量 \texttt{GeneralBikes}（$B^{\text{gen}}_{i,\tau}$）與電輔車數量 \texttt{ElectricBikes}（$B^{\text{elec}}_{i,\tau}$）。經空間多邊形邊界篩選後，臺大周邊研究場域共涵蓋 $\left|S\right| = 86$ 個唯一站點，2024 全年共計 $8{,}112{,}511$ 筆微觀狀態紀錄。

    \item \textbf{中央氣象署局屬氣象站地面觀測資料}：
    取自 TDX 平臺彙整之中央氣象署局屬氣象站觀測資料，2024 全年 366 個日檔案各包含全臺約 86 個測站、35 個氣象觀測欄位。本研究選取距臺大校總區約 $2.5$ 公里之局屬「臺北站」（\texttt{StationID = 466920}，位於中正區公園路）作為微氣候觀測基準，而非校內自動測站「臺灣大學站」（\texttt{A0A010}）。此一選擇係基於嚴格之資料覆蓋率檢核：臺灣大學站於 2024 年 1 月 1 日至 1 月 18 日完全無觀測紀錄，且於 2 月 9 日至 2 月 16 日存在連續缺測，全年僅涵蓋 340 天；反之，臺北站具備全年 366 天完整無間斷之 10 分鐘級觀測（共 $34{,}778$ 筆紀錄，每日中位數 $96$ 筆）。經交叉驗證兩站在重疊觀測期間之表現，兩站小時氣溫之皮爾森相關係數高達 $r = 0.993$（平均溫差僅 $0.4^\circ\text{C}$），日降雨事件分類一致率達 $95\%$，證實臺北站足以高度代表臺大生活圈之天氣動態。

    \item \textbf{學術行事曆與行政院人事行政總處辦公日曆表}：
    結合國立臺灣大學教務處公告之 112 學年度第 2 學期（2024-02-19 至 2024-06-09）與 113 學年度第 1 學期（2024-09-02 至 2024-12-20）行事曆，以及行政院人事行政總處公告之 2024 年政府行政機關辦公日曆表，建構研究場域專屬之日曆屬性表（\texttt{calendar\_2024.csv}，共 $366$ 日 $\times 5$ 欄位）。該表精確標記每日之星期代碼（\texttt{day\_of\_week}）、是否為實質上班上課日（\texttt{weekday}）、是否處於學期授課期間（\texttt{semester}）及放假或補班事由（\texttt{day\_off\_reason}）。2024 全年 366 天中，共包含 $244$ 個實質上班日（已扣除春節、清明、端午、中秋、國慶等國定假日與凱米、山陀兒、康芮三次颱風停班停課共 $19$ 個平日，並納入 2 月 17 日彈性補班日）與 $122$ 個非上班日；學期授課期間共 $222$ 天（含 $150$ 個學期中上班上課日），寒暑假期間則為 $144$ 天。
\end{enumerate}

\begin{table}[htbp]
\centering
\caption{研究採用之多源異質實證資料集規模與特徵彙整表}
\label{tab:data_sources_summary}
\small
\setlength{\tabcolsep}{2.0pt}
\begin{tabular}{@{} l c r r l @{}}
\toprule
\textbf{資料集名稱} & \textbf{期間} & \textbf{原始規模} & \textbf{研究區規模} & \textbf{時間解析度} \\
\midrule
1. YouBike 票證紀錄 & 2024--25 & $133{,}966{,}680$ 筆 & $86$ 站旅次 & 1 小時（歷時至秒） \\
\addlinespace[3pt]
2. TDX 車位快照 & 2024 年 & $111{,}802{,}383$ 筆 & $8{,}112{,}511$ 筆 & 變動觸發（中位 $1.9$ 分） \\
\addlinespace[3pt]
3. 氣象署臺北站觀測 & 2024 年 & $366$ 日全臺檔 & $34{,}778$ 筆 & 10 分（對齊 15/60 分） \\
\addlinespace[3pt]
4. 學術與行政行事曆 & 2024 年 & $366$ 日曆天 & $244$ 上班/$222$ 學期 & 1 日（含颱風與補班） \\
\bottomrule
\end{tabular}
\end{table}

\subsection{資料前處理與時空網格正則化管線}\label{subsec:data_preprocessing_pipeline}

由於政府開放資料在原始採集過程中存在採樣頻率改變、網路傳輸延遲、未記載狀態碼及特定日期缺檔等工程雜訊，若直接投入計量模型將引發嚴重之對齊偏誤。為此，本研究建置四階段可重現之資料前處理與時空網格正則化管線：

\begin{enumerate}
    \item \textbf{第一階段：研究場域空間拓撲篩選與實體欄位校正}
    \begin{enumerate}[label = \textbullet]
        \item \textbf{空間多邊形邊界定義與編碼檢核}：依據 WGS84 座標系（CRS 4326）定義由 $13$ 個頂點構成之封閉空間多邊形（涵蓋臺大校總區教學區、捷運公館站、捷運科技大樓站、長興街與基隆路男一/男六/男八舍宿舍群、溫州街生活圈及辛亥路/新生南路廊道），自全市站點中精確篩選出落於多邊形內之 $86$ 個唯一站點（因 3 個站點於年內微幅調整路側座標，於空間底圖中對應 $89$ 筆站點位置紀錄）。同時，針對票證刷卡資料中發現之 $218{,}870$ 筆（佔 $0.1634\%$）含 UTF-8 編碼異常（\texttt{U+FFFD} 替換字元）紀錄進行交叉比對，確認該批異常站名全數位於北投、士林等研究場域外，對本研究樣本造成零污染。
        \item \textbf{未記載營運狀態碼還原與時間戳記對齊}：TDX 原始快照之 \texttt{ServiceStatus} 欄位除官方文件定義之「正常營運」（$8{,}092{,}327$ 筆）外，出現未定義之數值代碼 \texttt{"2"}（共 $20{,}184$ 筆，涉及 $38$ 個站點）。經比對官方停用公告與站點車數分佈（代碼 \texttt{"2"} 期間平均可借車數仍達 $1.50$ 輛，並非單純空站），確認其為「停止營運（暫停服務）」並完成重編碼。此外，由於平臺接收時間 \texttt{UpdateTime} 相較於車柱設備感測時間 \texttt{SrcUpdateTime} 存在中位數 $24$ 秒之傳輸延遲（最大延遲達 $9{,}093$ 秒，且有 $4{,}187$ 筆出現時序倒置），所有時序排列與差分運算皆嚴格以 \texttt{SrcUpdateTime} 為唯一時間基準。
        \item \textbf{車種跨期口徑統一與動態容量識別}：檢核 TDX 快照發現，帳面可借總車數恆等式 $I_{i,\tau} = B^{\text{gen}}_{i,\tau} + B^{\text{elec}}_{i,\tau}$ 在 $8{,}112{,}511$ 筆紀錄中僅 $2$ 筆存在 $1$ 輛之瞬時誤差（一致率達 $99.99998\%$）。為與 2024 年前 10 個月未分車種之票證資料保持跨期口徑一致，本研究將一般車與電輔車合併計算帳面總庫存 $I_{i,\tau}$。同時，實證檢視發現同站點之「可借車數 $+$ 可還空柱數」（$I_{i,\tau} + O_{i,\tau}$）並非定值，而會隨個別車柱硬體鎖定或斷電而動態縮減，故站點物理總容量 $C_i$ 一律取自站點基本資料表（\texttt{stations\_in\_field.csv}），空柱狀態則直接採用快照實測之 $O_{i,\tau}$。
    \end{enumerate}

    \item \textbf{第二階段：變動觸發快照之時空網格正則化與三項防禦遮罩}
    \begin{enumerate}[label = \textbullet]
        \item \textbf{採樣機制結構性轉換與缺檔日盤點}：深度剖析 TDX 歷史資料發現，2024 年 1 月 1 日至 1 月 13 日期間，系統採「固定約 1 分鐘定時推播」機制（研究區每日約 $10$ 萬筆紀錄，最高達 $113{,}241$ 筆/日，1 月份單月即佔 $1{,}559{,}796$ 筆）；自 2024 年 1 月 14 日起，TDX 改採「狀態變動觸發記錄」機制，僅於站點可借車數或空柱數發生改變時寫入新快照（相鄰紀錄時間間隔中位數為 $1.9$ 分鐘，第 $90$ 百分位數為 $10.0$ 分鐘，研究區每日中位數降至 $21{,}028$ 筆）。此外，2024 全年 366 個原始日檔中，有 $10$ 天存在平臺端全區缺檔（含 10 月 31 日、11 月 2 日、11 月 7 日至 10 日共 $6$ 個 $0\text{ byte}$ 空檔，以及 11 月 1 日、11 月 3 日、11 月 11 日、12 月 10 日共 $4$ 個於研究場域內無紀錄之殘缺檔），實際有效觀測日為 $356$ 天。
        \item \textbf{平衡面板建立與前向填補防禦規則}：在變動觸發機制下，「無新快照寫入」在物理上即代表「站點車位狀態維持原狀」。基於此特性，本研究建立涵蓋 $86$ 個站點之雙尺度規則時空網格——包含小時級平衡面板（$86 \times 8{,}784 = 755{,}424$ 站時）與 15 分鐘級平衡面板（$86 \times 35{,}136 = 3{,}021{,}696$ 站時段）——並以各網格時點 $t$ 前最近一筆有效快照（$\text{\texttt{SrcUpdateTime}} \le t$）執行最後觀測值前向填補。為嚴防將缺檔日或設站前之無效狀態錯誤延續，管線強制實施三項不延續遮罩（將該網格設為 \texttt{NA}）：（1）\texttt{no record yet}：網格時點早於該站點首次出現觀測紀錄之時刻（如 2024-01-01 00:00 首筆快照為 00:01，或年內新設之 $6$ 個站點於啟用前之時段，最晚啟用者為 \texttt{TPE500106135} 於 10 月 15 日設站）；（2）\texttt{no-data day}：前向填補區間跨越上述 $10$ 個 TDX 平臺缺檔日之任何時段；（3）\texttt{stale > 24h}：距離最近一筆快照之狀態滯留時間（\texttt{status\_age\_min}）超過 $1{,}440$ 分鐘（24 小時）。經此嚴格過濾後，小時級與 15 分鐘級面板分別保留 $691{,}342$ 筆與 $2{,}765{,}475$ 筆有效觀測，有效覆蓋率皆達 $91.5\%$，且狀態滯留時間中位數僅 $6.8$ 分鐘。
    \end{enumerate}

    \item \textbf{第三階段：微觀步長差分與自然尖峰累積特徵檢驗}
    \begin{enumerate}[label = \textbullet]
        \item \textbf{微觀步長分佈與流量守恆檢驗}：對各站點依 \texttt{SrcUpdateTime} 排序後之 $8{,}112{,}425$ 個相鄰微觀步長計算單步庫存差分 $\Delta I_{i,\tau} = I_{i,\tau} - I_{i,\tau-1}$（嚴格排除跨缺檔日與間隔超過 24 小時之無效步長）。實證分佈顯示，高達 $99.9\%$ 之微觀步長滿足 $|\Delta I_{i,\tau}| \le 9$（$|\Delta I_{i,\tau}| \ge 10$ 之大步長跳躍僅佔全部步長之 $0.082\%$）。
        \item \textbf{尖峰高頻自然借還與卡車調度之實證辨析}：若以傳統文獻常用之經驗門檻 $\theta_{\text{rebal}} = 10$ 輛進行探索性拆解，在全年 $684{,}702$ 個正常營運有效站時中，僅 $5{,}537$ 個站時曾發生 $|\Delta I_{i,\tau}| \ge 10$ 之批次調度（全年累計調度移出入共 $85{,}975$ 輛，以捷運公館站 2 號出口 \texttt{TPE500101022} 之 $20{,}385$ 輛與 3 號出口 \texttt{TPE500101181} 之 $11{,}968$ 輛居冠）。更關鍵之實證發現為：在單小時淨流量絕對值達 $|\Delta I_{i,s}| \ge 25$ 輛之 $7{,}193$ 個劇烈變動站時中，僅 $18.5\%$（$1{,}330$ 個站時）包含卡車批次調度，其餘高達 $81.5\%$（$5{,}863$ 個站時）完全是由高頻率之自然借還（單小時中位數達 $33$ 次狀態更新、單步中位數變動 $6$ 輛）在換堂與通勤尖峰累積而成。此一實證特徵有力印證了第~\ref{sec:latent_demand_identification}~節之方法論主張：依賴固定庫存變動門檻切割調度與需求將嚴重誤傷尖峰自然流量，必須結合票證交易流水與局部波谷遮罩方能實現無偏識別。
    \end{enumerate}

    \item \textbf{第四階段：氣象異常碼清洗與跨日累積雨量校正}
    \begin{enumerate}[label = \textbullet]
        \item \textbf{異常代碼轉譯與雨跡識別}：針對中央氣象署臺北站（\texttt{466920}）之 $34{,}778$ 筆 10 分鐘觀測紀錄，首先將儀器檢修或通訊中斷代碼 \texttt{-99} 轉為缺漏值 \texttt{NA}（氣溫與相對濕度各 $54$ 筆、日累積雨量 $221$ 筆）。其次，針對日累積雨量欄位（\texttt{D24R}）出現之特殊代碼 \texttt{-990}（共 $5{,}998$ 筆，代表降水量低於 $0.1\text{ mm}$ 之微量「雨跡」），單獨建立布林指標 \texttt{D24R\_trace = TRUE} 並將數值歸零，避免負值代碼扭曲後續計數迴歸之降雨係數。
        \item \textbf{跨日歸零時間差校正與區間雨量還原}：由於氣象署自動雨量計係於每日 \texttt{00:10} 始執行跨日歸零，每日 \texttt{00:00} 之觀測讀數實為前一日之最終總累積雨量。管線透過將觀測時間減去 1 秒以對齊真實降雨歸屬日，再由單日累積雨量曲線之正差分還原出各 60 分鐘與 15 分鐘網格之真實區間降雨量（\texttt{precip\_mm}）與降雨二元指標（\texttt{is\_rainy}，2024 全年共計 $853$ 個降雨小時，佔比 $9.8\%$），最終透過時間鍵值與站點車位面板完成無縫接合。
    \end{enumerate}
\end{enumerate}

\section{無效庫存與調度事件識別}\label{sec:latent_demand_identification}

\subsection{潛在無效庫存與人為調度之識別困境}

在共享單車系統之實證研究中，欲由歷史營運資料還原使用者於各時段之期望需求率，核心原則在於必須以「站點實質能夠發生借還交互之有效暴露時間」作為計量分母。然而，站點庫存快照與自然借還交互之間存在兩項不可觀測之非隨機干擾：其一為站點車數因觸及某一下界而無法繼續向下借出之「單邊截斷底線」；其二為營運維運團隊以調度貨車或人工方式進行之「批次車輛搬移」。

在微觀時點 $\tau$（屬於巨觀決策時段 $s$），考量 YouBike 2.0 與 2.0E 電輔車共用相同之車柱與扣款介面，且 2024 年前 10 個月之票證資料未區分車種，站點 $i$ 回傳之帳面可借單車總數 $I_{i,\tau} = B^{\text{gen}}_{i,\tau} + B^{\text{elec}}_{i,\tau}$，在交互結構上可分解為「高於截斷底線之可交互車數 $B^{\text{eff}}_{i,\tau}$」與該時段內「無法完成向下借出交互之單邊截斷底線 $D_{i,s}$（涵蓋機械故障、低電量鎖定、車柱硬體受限或任何無法被借出之沉澱車數）」：
\begin{align}
I_{i,\tau} &= B^{\text{eff}}_{i,\tau} + D_{i,s}
\end{align}
（此處重點並非個別單車在機械上是否完好，而在於當車數降至 $D_{i,s}$ 時，站點向下的借車交互即發生單邊截斷。）與此同時，相鄰微觀時點間之帳面庫存變動量 $\Delta I_{i,\tau} = I_{i,\tau} - I_{i,\tau-1}$，則同時受自然借車量 $Y^+_{i,\tau}$、自然還車量 $Y^-_{i,\tau}$、調度補入量 $R^{\text{in}}_{i,\tau}$ 與調度移出量 $R^{\text{out}}_{i,\tau}$ 之交互驅動：
\begin{align}
\Delta I_{i,\tau} &= Y^-_{i,\tau} - Y^+_{i,\tau} + R^{\text{in}}_{i,\tau} - R^{\text{out}}_{i,\tau}
\end{align}

由於官方開放資料並未提供截斷底線明細與調度貨車之派工紀錄，$D_{i,s}$ 與 $(R^{\text{in}}_{i,\tau}, R^{\text{out}}_{i,\tau})$ 在計量模型中皆屬於不可觀測之潛在干擾變數。傳統文獻多仰賴研究者主觀設定之經驗啟發式門檻（例如假設車輛連續滯留超過特定時數即判定為故障車，或設定固定之單分鐘庫存跳躍門檻以切割調度事件）。然而，此類硬性推估極易在離峰時段將正常閒置車輛誤判為故障車，或在換堂尖峰將密集自然借還誤判為卡車調度，進而引入嚴重的測量誤差與內生性偏誤。

為徹底解決上述識別困境，本研究提出「基於局部波谷遮罩之需求識別策略」。本研究在理論架構中完整保留 $D_{i,s}$ 與 $(R^{\text{in}}_{i,\tau}, R^{\text{out}}_{i,\tau})$ 之物理定義，但透過結合「電子票證租借交易流水」與「庫存快照局部波谷下界」，在數學上嚴格證明真實需求率估計量對此二潛在干擾變數具備完全不變性，從而無須對干擾變數進行主觀經驗推估即可獲得零誤差之真實需求還原。

\subsection{局部庫存波谷與單邊截斷之包絡性質}\label{subsec:trough_envelope}

考慮站點 $i$ 於觀測時段 $s$ 內之微觀快照時間集合 $\mathcal{T}_s$。定義該時段內站點帳面可借總車數之局部庫存波谷為 $b_{i,s}$：
\begin{align}
b_{i,s} &\equiv \min_{\tau \in \mathcal{T}_s} I_{i,\tau}
\end{align}
依據可交互車數之非負性條件（$B^{\text{eff}}_{i,\tau} \ge 0$），站點在時段 $\mathcal{T}_s$ 內任意時點 $\tau$ 之帳面車數恆無法跌破該時段之單邊截斷底線 $D_{i,s}$：
\begin{align}
I_{i,\tau} &= B^{\text{eff}}_{i,\tau} + D_{i,s} \ge 0 + D_{i,s} = D_{i,s}, \quad \forall\ \tau \in \mathcal{T}_s
\end{align}
對全時段取最小值後，局部庫存波谷 $b_{i,s}$ 必然構成單邊截斷底線 $D_{i,s}$ 之上界包絡：
\begin{align}
b_{i,s} &= \min_{\tau \in \mathcal{T}_s} I_{i,\tau} \ge D_{i,s}
\end{align}
此一軌跡性質具有直接的物理意義：既然站點車數在時段 $s$ 內實際曾經向下減少並觸達 $b_{i,s}$，代表任何高於 $b_{i,s}$ 之水位皆未受單邊截斷阻擋；反之，當車數觸底停留於 $b_{i,s}$ 時，向下的借車交互即處於單邊截斷狀態（無論柱上剩餘車輛本身是否完好，該觸底期間皆無法繼續向下借出，理應自有效暴露分母中移除）。
基於此性質，本研究定義時段 $s$ 內站點 $i$ 之「真實可交互有效借車暴露時間集合」$\mathcal{T}^+_{i,s}$ 為帳面總庫存嚴格高於局部波谷之微觀時點集合：
\begin{align}
\mathcal{T}^+_{i,s} &\equiv \left\{ \tau \in \mathcal{T}_s \;\middle|\; I_{i,\tau} > b_{i,s} \right\}
\end{align}
當 $I_{i,\tau} > b_{i,s}$ 時，因車輛計數為離散整數，必有 $I_{i,\tau} \ge b_{i,s} + 1$。結合包絡不等式 $b_{i,s} \ge D_{i,s}$，即可確保站點在該時點高於截斷底線之可交互車數嚴格大於或等於 $1$：
\begin{align}
B^{\text{eff}}_{i,\tau} &= I_{i,\tau} - D_{i,s} \ge \left(b_{i,s} + 1\right) - b_{i,s} = 1
\end{align}
由於臺北市票證資料將借還時間戳記截斷為整點小時（僅能觀測時段內之借車總次數 $Y^+_{i,s}$，無從得知單一分鐘 $\tau$ 內是否發生借車），實務上直接由 TDX 快照序列統計滿足 $I_{i,\tau} > b_{i,s}$ 之分鐘數作為有效暴露時間 $\text{Exposure}^+_{i,s}$；在有效暴露時間集合 $\mathcal{T}^+_{i,s}$ 涵蓋之每一有效時點內，站點在物理上皆處於「未受單邊截斷、可實質進行借車交互」之狀態。

\subsection{現場營運交互狀況之穩健性檢驗}

基於庫存差分平衡式 $\Delta I_{i,\tau} = (Y^-_{i,\tau} - Y^+_{i,\tau}) + (R^{\text{in}}_{i,\tau} - R^{\text{out}}_{i,\tau})$ 與有效暴露集合 $\mathcal{T}^+_{i,s}$，以下檢驗五種現場營運狀況對需求率估計量 $\hat{\lambda}_{i,s} = Y^+_{i,s} / \text{Exposure}^+_{i,s}$ 之影響：

\begin{enumerate}
    \item \textbf{狀況一：維運移出與自然還車交錯（$R^{\text{out}}_{i,\tau} = 1, \, Y^-_{i,\tau} = 1$）}
    \begin{align}
    \Delta I_{i,\tau} &= \left(Y^-_{i,\tau} - Y^+_{i,\tau}\right) + \left(R^{\text{in}}_{i,\tau} - R^{\text{out}}_{i,\tau}\right) = (1 - 0) + (0 - 1) = 0, \qquad Y^+_{i,\tau} = 0
    \end{align}
    維運拉車（$-1$）與民眾還車（$+1$）一加一減相互抵銷，帳面車數不變（$I_{i,\tau} = I_{i,\tau-1}$）。還車僅計入 $Y^-_{i,s}$ 而不觸動借車分子 $Y^+_{i,s}$，分母暴露狀態亦維持原狀，故不影響 $\hat{\lambda}_{i,s}$。

    \item \textbf{狀況二：維運補入與自然借車交錯（$R^{\text{in}}_{i,\tau} = 1, \, Y^+_{i,\tau} = 1$）}
    \begin{align}
    \Delta I_{i,\tau} &= (0 - 1) + (1 - 0) = 0, \qquad I_{i,\tau} = I_{i,\tau-1} > b_{i,s} \;\Rightarrow\; \tau \in \mathcal{T}^+_{i,s}
    \end{align}
    供車狀態（$I_{i,\tau-1} > b_{i,s}$）下維運補車（$+1$）並於同分鐘被民眾借走（$-1$），快照車數雖因一加一減抵銷而無淨變動（$\Delta I_{i,\tau} = 0$），但因水位高於波谷故分母仍納入 $\mathcal{T}^+_{i,s}$（$+1$），且刷卡交易計入時段總分子 $Y^+_{i,s}$（$+1$），分子與分母同步對齊。

    \item \textbf{狀況三：故障與低電量無效車長期滯留（$B^{\text{eff}}_{i,\tau} = 0, \, I_{i,\tau} = D_{i,s}$）}
    \begin{align}
    I_{i,\tau} &= B^{\text{eff}}_{i,\tau} + D_{i,s} = D_{i,s} = b_{i,s} \;\Rightarrow\; \tau \notin \mathcal{T}^+_{i,s}
    \end{align}
    正常車耗盡後僅存 $D_{i,s}$ 輛沉澱車滯留，帳面車數收斂於波谷 $b_{i,s}$，波谷遮罩自動將該段期間自有效暴露時間 $\text{Exposure}^+_{i,s}$ 中剔除，不稀釋分母。

    \item \textbf{狀況四：貨車批次大量補車或移車（$R^{\text{in}}_{i,\tau} \gg 0$ 或 $R^{\text{out}}_{i,\tau} \gg 0$）}
    \begin{align}
    \Delta I_{i,\tau} &= R^{\text{in}}_{i,\tau} - R^{\text{out}}_{i,\tau}, \qquad Y^+_{i,\tau} = 0, \qquad \mathbb{I}\left(\tau \in \mathcal{T}^+_{i,s}\right) = \mathbb{I}\left(I_{i,\tau} > b_{i,s}\right)
    \end{align}
    即便調度作業橫跨數分鐘，工程搬運皆不經由民眾扣款系統，分子 $Y^+_{i,s}$ 不混入調度移出量；分母則隨快照車數 $I_{i,\tau}$ 即時反映站點有車可借之時長變化，故巨量調度跳躍不干擾單位時間借車率。

    \item \textbf{狀況五：自然尖峰需求借光導致空站截斷（$B^{\text{eff}}_{i,\tau} = 0 \Rightarrow I_{i,\tau} = b_{i,s}$）}
    \begin{align}
    \text{Exposure}^+_{i,s} &= \Delta t_s - \sum_{\tau \in \mathcal{T}_s} \mathbb{I}\left(I_{i,\tau} = b_{i,s}\right) = \Delta t_s - \text{EmptyMin}_{i,s}
    \end{align}
    尖峰需求將正常車借光後跌入波谷 $b_{i,s}$，斷貨期間之缺車分鐘數 $\text{EmptyMin}_{i,s}$ 自動自總時長 $\Delta t_s$ 中扣除，僅以斷貨前實質有車期間之借車數與暴露時間相除，消除截斷偏誤。
\end{enumerate}

表~\ref{tab:red_team_scenarios} 彙整五種營運狀況下分子、分母與估計量之反應：

\begin{table}[htbp]
\centering
\caption{現場營運狀況對分子、分母與真實需求率估計量之影響彙整表}
\label{tab:red_team_scenarios}
\small
\setlength{\tabcolsep}{3.2pt}
\begin{tabular}{@{} l l l c @{}}
\toprule
\textbf{營運狀況} & \textbf{分子 $Y^+_{i,s}$} & \textbf{分母 $\text{Exposure}^+_{i,s}$} & \textbf{偏誤 $\text{Bias}$} \\
\midrule
1. 維運移出與還車交錯 & 不受影響（還車寫入 $Y^-$） & 庫存淨差 $0$，維持原暴露 & $0$ \\
\addlinespace[3pt]
2. 維運補入與借車交錯 & 精確計入（刷卡 $+1$） & 高於波谷維持暴露（$+1$） & $0$ \\
\addlinespace[3pt]
3. 故障與沉澱車滯留 & 不受影響（無借車交易） & 波谷遮罩自動扣除無效期 & $0$ \\
\addlinespace[3pt]
4. 貨車批次大量調度 & 不受影響（工程無刷卡） & 快照即時反映供車時長 & $0$ \\
\addlinespace[3pt]
5. 自然尖峰空站截斷 & 完整累計有車期借車 & 自動扣除缺車分鐘數 & $0$ \\
\bottomrule
\end{tabular}
\end{table}

為具體說明上述五種現場營運狀況交錯發生時，模型如何自動分離出有效暴露時間，圖~\ref{fig:trough_exposure_example} 舉出一單一小時（$60$ 分鐘）之實務範例。設站內實際滯留 $2$ 輛未標記之故障車，期初帳面共 $6$ 輛車（$4$ 輛正常車、$2$ 輛故障車），全時段歷經三個階段：

\begin{figure}[!t]
\centering
\begin{tikzpicture}[
    xscale=0.188,
    yscale=0.39,
    >=Stealth,
    callout/.style={
        draw=#1!75!black,
        fill=#1!6!white,
        rounded corners=2pt,
        inner sep=2.2pt,
        font=\scriptsize,
        align=center,
        line width=0.5pt
    }
]
    % === 1. 背景分區（上層：安全供車區 vs 故障沉澱區） ===
    \fill[green!6!white] (0, 2) rectangle (60, 14.0);
    \fill[gray!22] (0, 0) rectangle (60, 2);

    % 階段垂直分隔虛線
    \draw[gray!55, dashed, line width=0.5pt] (15, -5.2) -- (15, 14.0);
    \draw[gray!55, dashed, line width=0.5pt] (30, -5.2) -- (30, 14.0);

    % 頂部階段標記
    \node[font=\scriptsize\bfseries, text=teal!65!black] at (7.5, 13.3) {階段 I：正常供車};
    \node[font=\scriptsize\bfseries, text=red!75!black] at (22.5, 13.3) {階段 II：斷貨觸底};
    \node[font=\scriptsize\bfseries, text=teal!65!black] at (45.0, 13.3) {階段 III：貨車調度與供車};

    % 水平網格線
    \foreach \y in {4, 6, 8, 10, 12} {
        \draw[gray!22, line width=0.3pt] (0, \y) -- (60, \y);
    }

    % 局部庫存波谷線
    \draw[red!80!black, dashed, line width=1.1pt] (0, 2) -- (60, 2);
    \node[fill=green!6!white, inner sep=1pt, font=\scriptsize\bfseries, text=red!80!black, anchor=south east] at (59.5, 2.15) {局部波谷線（$2$ 輛）};

    % 狀況三：底層故障車標記
    \node[font=\scriptsize\bfseries, text=gray!35!black] at (30, 1.0) {\textbf{狀況三}：$2$ 輛故障車長期滯留};

    % === 2. 帳面庫存軌跡 ===
    \draw[blue!75!black, line width=1.35pt]
        (0,6) -- (2,6) -- (2,5) -- (4.5,5) -- (4.5,4) -- (6,4) -- (6,5)
        -- (9,5) -- (9,4) -- (12,4) -- (12,3) -- (14,3) -- (14,4) -- (15,4) -- (15,2)
        -- (30,2)
        -- (30,12) -- (31.8,12) -- (31.8,10) -- (33.8,10) -- (33.8,9) -- (39,9) -- (39,8)
        -- (41,8) -- (41,9) -- (45,9) -- (45,8) -- (47,8) -- (47,9) -- (48,9)
        -- (48,5) -- (50,5) -- (50,4) -- (52,4) -- (52,5) -- (54,5) -- (54,4)
        -- (56,4) -- (56,5) -- (58,5) -- (58,4) -- (60,4);

    % 單行精簡標註（無算式、零重疊）
    % (1) 狀況一：tau = 8
    \filldraw[fill=teal!80!black, draw=white, line width=0.6pt] (8, 5) circle [x radius=0.45, y radius=0.22];
    \node[callout=teal, anchor=south west] (c1) at (1.5, 6.8) {\textbf{狀況一}：移車與還車抵銷};
    \draw[->, teal!75!black, line width=0.6pt] (8, 6.8) -- (8, 5.32);

    % (2) 狀況五：tau in (15, 30]
    \node[callout=red, anchor=south east] (c5) at (29.2, 9.8) {\textbf{狀況五}：好車借光觸底};
    \draw[->, red!75!black, line width=0.6pt] (28.0, 9.8) -- (28.0, 2.22);

    % (3) 狀況二：tau = 43 (供車期補車與借車抵銷)
    \filldraw[fill=orange!90!red, draw=white, line width=0.6pt] (43, 9) circle [x radius=0.45, y radius=0.22];
    \node[callout=orange, anchor=south] (c2) at (38.5, 11.3) {\textbf{狀況二}：補車與借車抵銷};
    \draw[->, orange!85!red, line width=0.6pt] (41.0, 11.3) -- (42.85, 9.25);

    % (4) 狀況四：tau = 30 補車 & tau = 48 移車
    \draw[->, purple!80!black, line width=1.6pt] (30, 2.2) -- (30, 11.8);
    \node[fill=green!6!white, inner sep=1pt, font=\scriptsize\bfseries, text=purple!80!black, anchor=west] at (30.4, 6.8) {$+10$};
    \draw[->, purple!80!black, line width=1.6pt] (48, 8.8) -- (48, 5.2);
    \node[fill=green!6!white, inner sep=1pt, font=\scriptsize\bfseries, text=purple!80!black, anchor=west] at (48.4, 7.0) {$-4$};
    \node[callout=purple] (c4) at (52.5, 11.3) {\textbf{狀況四}：貨車補車與移車};
    \draw[->, purple!75!black, line width=0.6pt] (c4.west) -- (30.4, 11.3);
    \draw[->, purple!75!black, line width=0.6pt] (48, 10.55) -- (48, 9.05);

    % 上層座標軸與刻度
    \draw[->, black, line width=0.8pt] (0, 0) -- (61.2, 0);
    \draw[->, black, line width=0.8pt] (0, -5.2) -- (0, 14.4);
    \node[font=\scriptsize\bfseries, anchor=south east] at (-0.2, 13.1) {帳面車數};
    \foreach \y in {0, 2, 4, 6, 8, 10, 12} {
        \draw[black, line width=0.6pt] (-0.4, \y) -- (0, \y);
        \node[font=\scriptsize, anchor=east] at (-0.45, \y) {$\y$};
    }

    % === 3. 中層：刷卡借車（分子：小時總計） ===
    \node[font=\scriptsize\bfseries, anchor=east] at (-0.5, -1.75) {刷卡借車};
    \node[font=\scriptsize, text=gray!70!black, anchor=east] at (-0.5, -2.55) {（分子）};
    \filldraw[fill=blue!8!white, draw=blue!65!black, line width=0.5pt] (0, -2.75) rectangle (60, -1.5);
    \node[font=\scriptsize\bfseries, text=blue!80!black] at (30.0, -2.12) {臺北市票證資料（僅記錄小時總量）：全時段累計刷卡借車 $Y^+_{i,s} = 18$ 次};

    % === 4. 下層：有效暴露（分母：TDX 快照序列） ===
    \node[font=\scriptsize\bfseries, anchor=east] at (-0.5, -3.8) {有效暴露};
    \node[font=\scriptsize, text=gray!70!black, anchor=east] at (-0.5, -4.6) {（分母）};

    % [0, 15] 有效暴露 (15 分)
    \filldraw[fill=teal!65!green!75, draw=teal!50!black, line width=0.5pt] (0, -4.7) rectangle (15, -3.45);
    \node[font=\scriptsize\bfseries, text=white] at (7.5, -4.08) {有效暴露 $15$ 分（$I > 2$）};

    % (15, 30] 遮蔽 (15 分)
    \filldraw[fill=gray!30, draw=gray!60!black, line width=0.5pt] (15, -4.7) rectangle (30, -3.45);
    \node[font=\scriptsize\bfseries, text=gray!30!black] at (22.5, -4.08) {波谷遮蔽 $15$ 分（$I = 2$）};

    % (30, 60] 有效暴露 (30 分)
    \filldraw[fill=teal!65!green!75, draw=teal!50!black, line width=0.5pt] (30, -4.7) rectangle (60, -3.45);
    \node[font=\scriptsize\bfseries, text=white] at (45.0, -4.08) {有效暴露 $30$ 分（$I > 2$）};

    % 時間軸刻度 (底部)
    \draw[->, black, line width=0.8pt] (0, -5.2) -- (61.2, -5.2);
    \node[font=\scriptsize\bfseries, anchor=west] at (61.3, -5.2) {時間（分）};
    \foreach \x in {0, 8, 15, 30, 43, 48, 60} {
        \draw[black, line width=0.6pt] (\x, -5.2) -- (\x, -5.55);
        \node[font=\scriptsize, anchor=north] at (\x, -5.55) {$\x$};
    }
\end{tikzpicture}
\caption{單一小時時段內五大現場營運狀況與有效借車暴露時間識別範例}
\label{fig:trough_exposure_example}
\end{figure}

\begin{enumerate}
    \item \textbf{階段 I（$0\text{--}15$ 分）：正常供車與維運抵銷（狀況一）}。第 $8$ 分鐘維運移出 $1$ 輛與民眾歸還 $1$ 輛相互抵銷，因均不觸發借車刷卡而不影響分子，且車數恆高於波谷線（$2$ 輛），全段 $15$ 分鐘計入有效暴露。

    \item \textbf{階段 II（$15\text{--}30$ 分）：好車借光觸底與故障車滯留（狀況三、五）}。第 $15$ 分鐘正常車借光，帳面僅剩 $2$ 輛故障車滯留並形成全時段波谷線（$b_{i,s} = 2$），系統直接依 $I_{i,\tau} = b_{i,s}$ 判定無車可借，自動遮蔽扣除全段 $15$ 分鐘缺車期。

    \item \textbf{階段 III（$30\text{--}60$ 分）：貨車批次調度與供車抵銷（狀況二、四）}。第 $30$ 與 $48$ 分鐘貨車分別批次補入 $10$ 輛與移出 $4$ 輛，工程搬運不產生刷卡紀錄故不干擾分子；第 $43$ 分鐘維運補車 $1$ 輛與民眾借車 $1$ 輛相互抵銷（狀況二），刷卡計入時段分子 $Y^+_{i,s}$，且全段車數維持於波谷線之上，完整計入 $30$ 分鐘有效暴露。
\end{enumerate}

綜觀全時段，臺北市票證資料記錄該小時刷卡借車共 $Y^+_{i,s} = 18$ 次。傳統做法因帳面車數恆大於零（$\ge 2$）而誤將 $60$ 分鐘全數計入分母，低估借車率為 $18 / 60 = 0.300$ 輛/分；本研究則由 TDX 快照自動扣除 $15$ 分鐘波谷缺車期，精確抓出 $15 + 30 = 45$ 分鐘有效暴露，還原真實需求率 $\hat{\lambda}_{i,s} = 18 / 45 = 0.400$ 輛/分。

\subsection{特殊邊界交互狀況之討論}\label{subsec:edge_case_interactions}

除上述五種常態營運狀況外，針對實務現場極少數可能發生之特殊邊界交互，本識別架構之處置機制與大樣本性質說明如下：

\begin{enumerate}
    \item \textbf{\textcolor{red}{同站異常試借之秒借秒還}}：若少數異常車尚未被鎖定，民眾刷出後始察覺車況不佳並立即推回同站歸還，帳面上將成對產生一筆借車（$Y^+ + 1$）與一筆還車（$Y^- + 1$）。本研究於資料前處理可直接過濾同站極短時試借紀錄；即便部分留存，成對發生之同站借還在站點淨流量（$\lambda_{i,s} - \mu_{i,s}$）中一加一減自然抵銷，於跨站轉移矩陣中亦僅為原地自迴圈（$i \to i$），完全維持系統流量守恆，且車柱具備短時歸還異常鎖定機制，於大樣本匯總下不影響估計結果。

    \item \textbf{單一時段內截斷底線之階梯跳變}：波谷遮罩係假設單邊截斷底線 $D_{i,s}$ 於單一決策時段 $s$（15 或 60 分鐘）內維持局部平穩；若同一時段內先於低水位觸底，中途因壞車還入（或維運清運壞車）使截斷底線墊高（或降低）並再次觸底，單一水平波谷 $b_{i,s}$ 僅能捕捉較低之截斷水位。然而，車輛損壞與維修清運本身皆屬低頻率事件，在單一短時段內同時發生「先觸底、中途底線跳變、再二次觸底」之複合機率極低，於跨期大樣本匯總下屬可忽略之微小擾動。

    \item \textbf{還車端車柱鎖定與分鐘網格之微觀時間差}：在還車端，理論上雖可定義對稱之局部空柱波谷 $o_{i,s} \equiv \min_{\tau \in \mathcal{T}_s} O_{i,\tau}$，惟實務上車柱斷電或硬體鎖定時，TDX 回傳之可還空柱數 $O_{i,\tau}$ 已由系統自動即時扣減（反映於第~\ref{subsec:data_preprocessing_pipeline}~節之 $I_{i,\tau} + O_{i,\tau} < C_i$，與故障車仍停於柱上計入 $I_{i,\tau}$ 之機制不同），且多數未滿站時段之 $o_{i,s} > 0$ 係源於充裕空柱而非故障卡死；故直接以 $O_{i,\tau} = 0$ 即可避免過度遮蔽充裕空柱時段並捕捉實質滿站截斷。至於極少數發生於波谷（$I_{i,\tau} = b_{i,s}$）或滿站（$O_{i,\tau} = 0$）交界分鐘內之「一進一出抵銷」，雖因票證資料僅記錄小時總量而無從逐分識別，惟此類交界瞬間於單小時內至多發生 1～2 次，於全年大樣本網格匯總下皆自然平滑抵銷。
\end{enumerate}

\section{截斷借還需求之還原}\label{sec:demand_reconstruction}

\subsection{基準抵達過程之設定動機與資料物理邊界}\label{subsec:poisson_data_boundary}

在建立潛在借還需求之統計證明前，須先由雙源資料之資訊集邊界出發，闡明本研究以「分段常數齊次 Poisson 過程」建模之計量必然性，以及其與機器學習預測模型之分工關係：

\begin{enumerate}
    \item \textbf{次時段時變強度之不可辨識性與最大熵基準}：
    微觀交通流之抵達強度雖隨時間波動（非齊次強度 $\lambda_i(t)$），惟配適連續時變函數 $\lambda_i(t)$ 必須具備分秒級之微觀交易時間戳記。受限於官方票證資料將時間戳記截斷至整點（\texttt{:00:00}），且 TDX 快照差分 $\Delta I_{i,\tau}$ 混雜未記載之卡車調度 $(R^{\text{in}}_{i,\tau}, R^{\text{out}}_{i,\tau})$ 與截斷底線 $D_{i,s}$，實證上在有效暴露時間集合 $\mathcal{T}^+_{i,s}$ 內僅能觀測到時段交易總量 $Y^+_{i,s}$，無法取得次時段內之微觀事件分佈。在「僅具區間總量與暴露時長」之資訊集約束下，小時內之連續波形 $\lambda_i(t)$ 在統計上屬不可辨識。依據最大熵原理，將單一決策時段 $s$ 內之潛在抵達強度設為局部平穩常數 $\lambda_{i,s}$（分段常數齊次 Poisson 過程），是不引入武斷函數假設之最簡無偏基準，亦直接對應調度車隊以小時為週期之作業尺度。

    \item \textbf{目標標籤去偏與機器學習之上下游分工}：
    監督式機器學習必須建立於未受污染之目標標籤之上。若直接以原始票證計數 $Y^+_{i,s}$ 訓練模型，將把斷貨截斷導致之零交易誤判為無需求；若直接擬合 TDX 庫存變動 $\Delta I_{i,\tau}$，則會將卡車調度搬移誤判為自然借還人潮。即便先以波谷遮罩剔除截斷時段，各站剩餘之有效暴露時長並不相等（$\text{Exposure}^+_{i,s} \le \Delta t_s$）且僅具單一總量，機器學習無法繞過資料限制推演出微觀時變率，仍須先經由暴露時長標準化以消除截面偏誤。因此，本節推導與第~\ref{chap:Dynamic Rebalancing Models & Algorithms}~章之預測模型屬嚴格之上下游關係：本節藉由點過程理論剝離硬體截斷與調度干擾，產出無偏之小時級真實需求率標籤；後續計量與機器學習模型方能在此純淨標籤上，進階學習氣象、日曆與空間拓撲之非線性關係。
\end{enumerate}

\subsection{真實需求率之無偏性與干擾參數不變性證明}

為將上述五種營運狀況之零誤差特性予以一般化，本節建立嚴格之計量統計證明，推導暴露量校正需求率估計量之無偏性與對未知干擾參數之不變性。

設站點 $i$ 於時段 $s$（總時長 $\Delta t_s$）內，使用者之潛在無約束借車抵達行為外生於車柱硬體狀態，服從強度為 $\lambda_{i,s}$ 之齊次 Poisson 點過程 $\{N_{i,s}(t), t \in [0, \Delta t_s]\}$。給定時點 $t^-$ 前之歷史資訊過濾集 $\mathcal{F}_{t^-}$，由獨立增量性質，無窮小時間增量 $dt$ 內之條件期望抵達數滿足：
\begin{align}
E\left[dN_{i,s}(t) \;\middle|\; \mathcal{F}_{t^-}\right] &= \lambda_{i,s} \, dt
\end{align}
令 $dY^+_{i,s}(t)$ 為電子票證租借資料庫在時點 $t$ 實際記錄到之借車交易增量。由於維運人員之貨車調度作業 $(R^{\text{in}}_{i,\tau}, R^{\text{out}}_{i,\tau})$ 不產生票證刷卡紀錄，且潛在使用者抵達站點時，唯有當站點處於未受單邊截斷之可交互狀態（即 $B^{\text{eff}}_{i,t^-} = I_{i,t^-} - D_{i,s} \ge 1$）時方能實際完成借車交互，故觀測借車過程與潛在抵達過程之結構關係式可嚴格寫為：
\begin{align}
dY^+_{i,s}(t) &= \mathbb{I}\left(I_{i,t^-} - D_{i,s} \ge 1\right) \cdot dN_{i,s}(t)
\end{align}
其中 $\mathbb{I}(\cdot)$ 為指示函數。

現透過波谷遮罩將觀測時間軸限制於有效暴露時間集合 $\mathcal{T}^+_{i,s}$。由第~\ref{subsec:trough_envelope}~節之單邊截斷包絡性質 $b_{i,s} \ge D_{i,s}$ 可知，$\forall\ t \in \mathcal{T}^+_{i,s}$（忽略單分鐘交界瞬間之高階網格離散化微差），站點於借車發生前夕之左連續可交互車數恆滿足 $I_{i,t^-} - D_{i,s} \ge 1$，是以 $\mathcal{F}_{t^-}$-可預測之指示函數在 $\mathcal{T}^+_{i,s}$ 上恆等於 $1$：
\begin{align}
\mathbb{I}\left(I_{i,t^-} - D_{i,s} \ge 1\right) &= 1, \quad \forall\ t \in \mathcal{T}^+_{i,s}
\end{align}
對有效暴露集合 $\mathcal{T}^+_{i,s}$ 進行積分，即可得時段 $s$ 內之觀測借車總次數 $Y^+_{i,s}$：
\begin{align}
Y^+_{i,s} &= \int_{t \in \mathcal{T}^+_{i,s}} dY^+_{i,s}(t) = \int_{t \in \mathcal{T}^+_{i,s}} 1 \cdot dN_{i,s}(t) = \int_{t \in \mathcal{T}^+_{i,s}} dN_{i,s}(t)
\end{align}
由計數過程之 Doob-Meyer 分解與獨立增量性質，即便於未完全斷貨時段 $b_{i,s} > D_{i,s}$ 對最低水位產生保守修剪，只要有效暴露區間內指示函數恆為 $1$，子區間上之抵達強度仍保持 $\lambda_{i,s}$ 不變。給定有效暴露時間集合 $\mathcal{T}^+_{i,s}$（其總測度即為有效暴露分鐘數 $\text{Exposure}^+_{i,s} = \int_{t \in \mathcal{T}^+_{i,s}} 1 \, dt = \Delta t_s - \text{EmptyMin}_{i,s}$）、單邊截斷底線 $D_{i,s}$ 與調度跳躍量 $(R^{\text{in}}_{i,s}, R^{\text{out}}_{i,s})$，求取 $Y^+_{i,s}$ 之條件期望值可得：
\begin{equation}
\begin{aligned}
E\left[Y^+_{i,s} \;\middle|\; \mathcal{T}^+_{i,s}, D_{i,s}, R^{\text{in}}_{i,s}, R^{\text{out}}_{i,s}\right] &= \int_{t \in \mathcal{T}^+_{i,s}} E\left[dN_{i,s}(t) \mid \mathcal{F}_{t^-}\right] = \int_{t \in \mathcal{T}^+_{i,s}} \lambda_{i,s} \, dt \\
&= \lambda_{i,s} \cdot \text{Exposure}^+_{i,s}
\end{aligned}
\end{equation}
由此定義站點 $i$ 於時段 $s$ 之真實潛在借車率估計量 $\hat{\lambda}_{i,s}$ 為：
\begin{align}
\hat{\lambda}_{i,s} &= \frac{Y^+_{i,s}}{\text{Exposure}^+_{i,s}}
\end{align}
當 $\text{Exposure}^+_{i,s} > 0$ 時，在給定有效暴露時間 $\mathcal{T}^+_{i,s}$（或於跨期大樣本匯總 $\sum Y^+_{i,s} / \sum \text{Exposure}^+_{i,s}$ 下由線性期望與大數法則求取極限）對該估計量取期望值，即可得無偏性與大樣本一致性：
\begin{equation}
\begin{aligned}
E\left[\hat{\lambda}_{i,s} \;\middle|\; \mathcal{T}^+_{i,s}, D_{i,s}, R^{\text{in}}_{i,s}, R^{\text{out}}_{i,s}\right] &= \frac{E\left[Y^+_{i,s} \;\middle|\; \mathcal{T}^+_{i,s}, D_{i,s}, R^{\text{in}}_{i,s}, R^{\text{out}}_{i,s}\right]}{\text{Exposure}^+_{i,s}} \\
&= \frac{\lambda_{i,s} \cdot \text{Exposure}^+_{i,s}}{\text{Exposure}^+_{i,s}} = \lambda_{i,s}
\end{aligned}
\end{equation}
進而計算估計量之統計偏誤：
\begin{align}
\text{Bias}\left(\hat{\lambda}_{i,s}\right) &= E\left[\hat{\lambda}_{i,s}\right] - \lambda_{i,s} = \lambda_{i,s} - \lambda_{i,s} = 0
\end{align}
同時，對不可觀測之單邊截斷底線 $D_{i,s}$ 及人為調度跳躍量 $R^{\text{in}}_{i,s}, R^{\text{out}}_{i,s}$ 取偏導數，可得：
\begin{align}
\frac{\partial E\left[\hat{\lambda}_{i,s}\right]}{\partial D_{i,s}} &= 0, \qquad \frac{\partial E\left[\hat{\lambda}_{i,s}\right]}{\partial R^{\text{in}}_{i,s}} = 0, \qquad \frac{\partial E\left[\hat{\lambda}_{i,s}\right]}{\partial R^{\text{out}}_{i,s}} = 0
\end{align}
此式表明，潛在干擾變數 $\left(D_{i,s}, R^{\text{in}}_{i,s}, R^{\text{out}}_{i,s}\right)$ 僅透過改變分母有效暴露時長 $\text{Exposure}^+_{i,s}$ 影響估計量之有限樣本變異數（$\text{Var}\left(\hat{\lambda}_{i,s} \mid \mathcal{T}^+_{i,s}\right) = \lambda_{i,s} / \text{Exposure}^+_{i,s}$），而與估計量之一階期望值嚴格正交，故不引入系統性截斷偏誤。

同理，在還車需求端，使用者歸還單車僅受限於站點是否處於滿站單邊截斷狀態（無空柱可還）。對稱定義有效還車暴露時間集合 $\mathcal{T}^-_{i,s} \equiv \{\tau \in \mathcal{T}_s \mid O_{i,\tau} > 0\}$，其總測度為有效還車暴露分鐘數 $\text{Exposure}^-_{i,s} = \Delta t_s - \text{FullMin}_{i,s}$（其中 $\text{FullMin}_{i,s}$ 為時段 $s$ 內 $O_{i,\tau} = 0$ 之滿站截斷分鐘數），並以票證資料庫記錄之實際還車次數 $Y^-_{i,s}$ 除以有效還車暴露時間，即可對稱獲得真實潛在還車率估計量：
\begin{align}
\hat{\mu}_{i,s} &= \frac{Y^-_{i,s}}{\text{Exposure}^-_{i,s}}, \qquad E\left[\hat{\mu}_{i,s}\right] = \mu_{i,s}
\end{align}

由上述嚴格證明，我們可以得出以下核心結論：無論在借車端抑或還車端，只要將觀測時間軸限制於站點實質具備交互能力之有效暴露集合（未受借車截斷之 $\mathcal{T}^+_{i,s}$ 與未受還車截斷之 $\mathcal{T}^-_{i,s}$），估計量即具備嚴格之統計無偏性與一致性：
\begin{empheq}[box=\fbox]{align}
E\left[\hat{\lambda}_{i,s}\right] = \lambda_{i,s}, \qquad E\left[\hat{\mu}_{i,s}\right] = \mu_{i,s}
\end{empheq}
更為本質的是，估計量之期望值對系統中不可觀測之單邊截斷底線 $D_{i,s}$ 及人為貨車調度量 $(R^{\text{in}}_{i,s}, R^{\text{out}}_{i,s})$ 之偏導數恆等於零。此一數學性質在理論與實務上確立了以下特性：
\begin{enumerate}
    \item \textbf{干擾參數完全不變性}：估計量之期望值嚴格獨立於截斷底線與卡車調度之存在與否，證明本方法在數學結構上不受現場營運雜訊影響，無須依賴文獻常見之主觀經驗門檻進行人為修剪。
    \item \textbf{微觀對帳零偏誤}：藉由波谷包絡與雙資料通道解耦，成功將工程搬運與自然交易進行物理隔離，使尖峰時段借車分子 $Y^+_{i,s}$ 與分母有效暴露時間 $\text{Exposure}^+_{i,s}$ 精確對齊，徹底消除因觸底單邊截斷所造成的估計偏誤。
    \item \textbf{實證可計算性}：實證研究者無須耗費龐大算力去「猜測」個別單車是否損壞或追蹤調度車輛軌跡，僅憑公開之刷卡流水與站點快照即可直接計算出真實無偏之借還需求率。\hfill$\blacksquare$
\end{enumerate}

\subsection{實證估算之最小充分統計量與干擾參數消除}

依據機率統計中之 Neyman-Fisher 因子分解定理，在給定有效暴露集合 $\mathcal{T}^+_{i,s}$ 下，由於觀測借車次數服從參數為 $\lambda_{i,s} \cdot \text{Exposure}^+_{i,s}$ 之計數分佈，其對數概似函數僅透過觀測交易總數 $Y^+_{i,s}$ 與有效暴露時間 $\text{Exposure}^+_{i,s}$ 依賴於樣本資料：
\begin{equation}
\begin{aligned}
\ln \mathcal{L}\left(\lambda_{i,s} \;\middle|\; Y^+_{i,s}, \text{Exposure}^+_{i,s}, D_{i,s}, R^{\text{in}}_{i,s}, R^{\text{out}}_{i,s}\right) &= Y^+_{i,s} \ln\left(\lambda_{i,s} \cdot \text{Exposure}^+_{i,s}\right) \\
&\quad - \lambda_{i,s} \cdot \text{Exposure}^+_{i,s} - \ln\left(Y^+_{i,s}!\right)
\end{aligned}
\end{equation}
此一數學結構嚴格證明：二元數據配對 $\left(Y^+_{i,s}, \text{Exposure}^+_{i,s}\right)$ 與 $\left(Y^-_{i,s}, \text{Exposure}^-_{i,s}\right)$ 即為識別真實借還需求率 $\left(\lambda_{i,s}, \mu_{i,s}\right)$ 之「最小充分統計量」。

基於上述不變性與充分性證明，本研究在實證資料處理管線中，實現了對不可觀測干擾參數之完全消除：
\begin{enumerate}
    \item \textbf{實證唯一必算之核心充分數據}：在構建各站點各時段之需求分析矩陣時，針對單站供需狀態，僅需由電子票證租借流水帳匯總實際借車次數 $Y^+_{i,s}$ 與還車次數 $Y^-_{i,s}$，並由 TDX 微觀庫存快照提取高於局部波谷之有效借車暴露分鐘數 $\text{Exposure}^+_{i,s}$ 與空柱大於零之有效還車暴露分鐘數 $\text{Exposure}^-_{i,s}$。
    \item \textbf{理論保留但實證免算之干擾變數}：單邊截斷底線 $D_{i,s}$ 與調度貨車補移車跳躍量 $\left(R^{\text{in}}_{i,s}, R^{\text{out}}_{i,s}\right)$ 雖於理論模型中保留其符號以完整刻畫物理系統，但在求算真實需求率時，因其在概似函數與估計式中之偏導數恆為零，實證上完全無須編寫啟發式規則進行推估或剝離。
\end{enumerate}
此一設計不僅大幅簡化了巨量歷史資料之前處理運算複雜度，更重要的是從根源上杜絕了人為設定過濾門檻所造成的估計偏誤，為第~\ref{chap:Dynamic Rebalancing Models & Algorithms}~章之時變計數迴歸與非齊次生滅佇列模型提供統計上嚴格無偏之真實需求輸入。
